\documentclass[10pt,a4paper]{amsart}
\usepackage[margin=19mm]{geometry}
\usepackage{amsmath,amssymb,mathtools}
\usepackage[scr=rsfs,cal=euler]{mathalfa}
\usepackage{xfrac}
\usepackage[unicode,hidelinks,bookmarks=false]{hyperref}
\newcommand{\ie}{i.e.\ }
\newcommand{\cf}{cf.\ }
\newcommand{\resp}{respectively\ }
\newcommand{\Subsection}{\subsection}
\providecommand{\nobreakdash}{\nobreak\hbox{-}}
\setlength{\emergencystretch}{2em}
\multlinegap0pt
\usepackage{bm}
\usepackage{bbm}
\usepackage{dsfont}
\usepackage{xspace}
\usepackage{cancel}
\usepackage{mathtools}
\usepackage{amsthm}
\makeatletter
\@ifundefined{definition}{\newtheorem{definition}{Definition}}{}
\@ifundefined{proposition}{\newtheorem{proposition}{Proposition}}{}
\makeatother
\title[Poincar\'e duality for singular tropical hypersurfaces]{Poincar\'e duality for singular tropical hypersurfaces}
\author{Samuel Dentan}
\date{Source: arXiv:2512.24548v2 (CC BY 4.0)}
\begin{document}
\maketitle

\noindent\emph{Source and licence.} Poincar\'e duality for singular tropical hypersurfaces. Version arXiv:2512.24548v2 (CC BY 4.0), distributed under the Creative Commons Attribution 4.0 International licence, as recorded in the arXiv licence metadata for this exact version and shown on the abstract page https://arxiv.org/abs/2512.24548v2. Full rights notice: licenses/arxiv-2512.24548-CC-BY-4.0.txt.\par\smallskip

\noindent\textit{Editorial scope.} Counted unit: the Proposition whose source label is \texttt{ranphi} in the article (source0.tex lines 1658--1665, source label \texttt{ranphi}), delivered together with the article's own complete original proof of it (source0.tex lines 1699--1778, one \texttt{proof} environment of 80 lines). No mathematics is written, repaired or re-derived: statement and proof are the article's own text line for line. The proof is detached from the statement in the source: the article states the result at lines 1658--1665 and prints its proof at lines 1699--1778, and the proof environment's own header names the statement, so the association is the article's own. Required context retained uncounted: max (lines 740--747); defphi (lines 1126--1141); sat (lines 1349--1355). Every definition, display and label of those lines is retained unchanged. Omitted from this package and not redistributed: 1--739, 748--1125, 1142--1348, 1356--1657, 1666--1698, 1779--1973. Nothing inside a retained excerpt is omitted or altered: every retained body is delivered exactly as the article's lines are written, so no omission receipt of any kind accompanies this package. Citations: the retained excerpts carry no citation command at all, so no bibliography entry of the article is transcribed and no reference is redistributed. Reference adaptation: none was needed and none was made; every reference inside the delivered unit resolves inside it, so no unresolved reference or citation is produced. Printed slips kept literal: no printed word, symbol, hypothesis or number of the counted statement or of its proof was altered. Layout and class: the article's own class is \texttt{article} with packages including geometry, inputenc, amsfonts, graphicx, placeins, float, subfigure, amsmath, amsthm, amssymb, adjustbox, xcolor, ulem, tikz-cd; the wrapper is the lane's pinned \texttt{amsart} editorial wrapper at 10pt on A4, which restates the theorem-like environments the retained text needs and adds the article's own macro definitions that the retained text needs (none), with extra wrapper packages (bm, bbm, dsfont, xspace, cancel, mathtools). The delivered numbering is the unit's own. Assets: no retained span contains a figure environment, a graphics-inclusion command, a PDF-page inclusion, a table or a TikZ picture; no image file is copied into this package, the package declares no asset dependency and no image file is redistributed with it. Licence: version arXiv:2512.24548v2 of this article is distributed under the Creative Commons Attribution 4.0 International licence, as recorded in the arXiv licence metadata for this exact version and shown on the abstract page https://arxiv.org/abs/2512.24548v2. A full rights notice is delivered at licenses/arxiv-2512.24548-CC-BY-4.0.txt. Characters: every retained excerpt is pure ASCII, so the delivered engine, XeLaTeX with its default Latin Modern text font, reports no missing character.
\begin{definition}
    (Hypothesis 1 and 2)
    
    \underline{Hypothesis 1:} We say that the triangulation $\Gamma$ satisfies the hypothesis 1 for the integral domain $R$ if for any edge $\sigma^1 \in \Gamma$ the reduced length $|\sigma^1|_r$ is non-zero in the integral domain $R$.
    
    \underline{Hypothesis 2:} We say that the triangulation $\Gamma$ satisfies the hypothesis 2 for the integral domain $R$ if the reduced length $|\sigma^1|_r$ is invertible in $R$.
    \label{max}
\end{definition}

\begin{definition}
    (Cap-product with $[\Omega]$) For any $n$-dimensional polytope $P$ and triangulation $\Gamma$ with integral vertices we define the following linear maps:
    \begin{equation*}
        \forall \sigma^{n-q}\in \Gamma,~~ \phi_{\sigma^{n-q}}: \begin{array}{rcl}
             C^{q}(\mathcal{F}^{p}_R(\sigma^{n-q},*)) & \rightarrow & C_{n-1-q}(\mathcal{F}_{n-1-p}^R(*,\sigma^{n-q}))  \\
             \beta & \mapsto & \beta \cap [\Omega] ~;
        \end{array}
    \end{equation*}
    \begin{equation*}
        \phi^q: \begin{array}{rcl}
             C^{q}\mathcal{F}^{p}_R) & \rightarrow & C_{n-1-q}(\mathcal{F}_{n-1-p}^R)  \\
             \beta & \mapsto & \beta \cap [\Omega] 
        \end{array}.
    \end{equation*}
    \label{defphi}
\end{definition}

\begin{definition}
    (Saturation) If $R$ is an integral domain, if $L$ is a $R$-module, and $S$ a sub-module of $L$, we say that $S$ is saturated in $L$ if:
    \begin{equation*}
        \forall y\in L, \forall \lambda\in R \backslash\{0\},~~\lambda y \in S \implies y \in S.
    \end{equation*}
    \label{sat}
\end{definition}

\begin{proposition}
    Let us fix a simplex $\sigma^a \in \Gamma$. Assume that the triangulation satisfies the hypothesis 1 for the integral domain $R$ (see Definition \ref{max}). Let us denote $V=\sum_{\sigma^1 \subseteq \sigma^a}T_R \sigma^1$. Then the map $\phi_{\sigma^{a},(p)}$ (see Definition \ref{defphi}) satisfies:
    \begin{equation*}
        \mathrm{rank} ~ \phi_{\sigma^{a},(p)} =\binom{\mathrm{dim}F_{\sigma^a}}{p-n+\mathrm{dim}F_{\sigma^a}}-\binom{\mathrm{dim}F_{\sigma^a}-\mathrm{dim}V}{p-n+\mathrm{dim}F_{\sigma^a}-\mathrm{dim}V}.
    \end{equation*}
    Moreover, if we assume that the triangulation satisfies the hypothesis 2 for the integral domain $R$ (see Definition \ref{max}), the image of $\phi_{\sigma^{a},(p)}$ is a $R$-module saturated in $C_{a-1}(\mathcal{F}_{n-1-p}^R(*,\sigma^a))$ (see Definition \ref{sat} for the saturation of a sub-module).
    \label{ranphi}
\end{proposition}

\begin{proof}
    Proof of Proposition \ref{ranphi}.
    
    Let us fix a basis $(f_1,...,f_{\mathrm{rank} V})$ of $V$, and let $\{f_{\mathrm{rank}V+1},...,f_{\mathrm{dim}F_{\sigma^a}}\}$ a set of vectors which complete it, to obtain a basis $(f_1,...,f_{\mathrm{dim} F_{\sigma^a}})$ of $T_R F_{\sigma^a}$. Let $\{e_1,...,e_{n-\mathrm{dim}F_{\sigma^a}}\}$ a set of vectors which complete $(f_1,...,f_{\mathrm{dim} F_{\sigma^a}})$, to obtain a basis of $M_R$. Let us denote $(f_1^*,...,f_{\mathrm{dim} F_{\sigma^a}}^*,e_1^*,...,e_{n-\mathrm{dim} F_{\sigma^a}}^*)$ the dual basis of $N_R$. For every $\sigma^1 \subseteq \sigma^a$, let us denote $u_{\sigma^1}$ a generator of $T_R\sigma^1$. We can complete it to obtain a basis $(u_{\sigma^1},v_1^{\sigma^1},...,v_{\mathrm{rank}V -1}^{\sigma^1})$ of $V$, and then the following basis of $M_R$:
    \begin{equation*}
        (u_{\sigma^1},v_1^{\sigma^1},...,v_{\mathrm{rank}V -1}^{\sigma^1},f_{\mathrm{rank}V+1},...,f_{\mathrm{dim}F_{\sigma^a}},e_1,...e_{n-\mathrm{dim}F_{\sigma^a}}).
    \end{equation*}
    The dual basis is of the form:
    \begin{equation*}
        (u_{\sigma^1}^*,(v_1^{\sigma^1})^*,...,(v_{\mathrm{rank}V -1}^{\sigma^1})^*,f_{\mathrm{rank}V+1}^*,...,f_{\mathrm{dim}F_{\sigma^a}}^*,e_1^*,...e_{n-\mathrm{rank}V}^*)
    \end{equation*}
    where the vectors $f_{\mathrm{rank}V+1}^*,...,f_{\mathrm{dim}F_{\sigma^a}}^*,e_1^*,...e_{n-\mathrm{dim}F_{\sigma^a}}^*$ are the same as the ones that appear in the dual basis $(f_1^*,...,f_{\mathrm{dim} F_{\sigma^a}}^*,e_1^*,...,e_{n-\mathrm{dim} F_{\sigma^a}}^*)$.
    
    Then $((v_1^{\sigma^1})^*, ..., (v_{\mathrm{dim}F_{\sigma^a} -1}^{\sigma^1})^*, f_{\mathrm{rank}V+1}^*, ..., f_{\mathrm{dim}F_{\sigma^a}}^*, e_1^* ,..., e_{n-\mathrm{dim}F_{\sigma^a}}^*)$ is a basis of $T_R \sigma^{1\perp}$. Then $(v_1^{\sigma^1})^*\wedge...\wedge(v_{\mathrm{dim}F_{\sigma^a}-1}^{\sigma^1})^*\wedge f_{\mathrm{rank}V+1}^*\wedge...\wedge f_{\mathrm{dim}F_{\sigma^a}} ^*\wedge e_1^*\wedge... \wedge e_{n-\mathrm{dim}F_{\sigma^a}}^*$ and $g((o(\sigma^1)) \otimes 1_R$ are both generators of $\bigwedge^{n-1}T_R\sigma^{1\perp}$, thus there exists an invertible element $\lambda_{\sigma^1}$ of $R$ such that:
    \begin{align*}
        &(v_1^{\sigma^1})^*\wedge...\wedge(v_{\mathrm{dim}F_{\sigma^a} -1}^{\sigma^1})^*\wedge f_{\mathrm{rank}V+1}^* \wedge...\wedge f_{\mathrm{dim}F_{\sigma^a}}^* \wedge e_1^*\wedge... \wedge e_{n-\mathrm{dim}F_{\sigma^a}}^* =\lambda_{\sigma^1} g(o(\sigma^1)) \otimes 1_R.
    \end{align*}

    Let $\alpha \in \bigwedge^p M_R$. We can write:
    \begin{equation*}
        \alpha = \sum_{k=0}^p\sum_{A \in \mathcal{P}_{p-k}(\{1,...,\mathrm{dim}F_{\sigma^a}\})}\sum_{B \in \mathcal{P}_{k}(\{1,...,n-\mathrm{dim} F_{\sigma^a}\})}\alpha_{A,B} \bigwedge_{i \in A}f_i \wedge\bigwedge_{j\in B} e_i
    \end{equation*}
    where $\mathcal{P}_l(\{1,...,s\})$ denotes the set of subsets of cardinal $l$ of $\{1,...,s\}$. However, for every $k\neq n- \mathrm{dim} F_{\sigma^a}$, $B \in \mathcal{P}_{k}(\{1,...,n-\mathrm{dim} F_{\sigma^a}\})$, and $A \in \mathcal{P}_{p-k}(\{1,...,\mathrm{dim}F_{\sigma^a}\})$, for every $\sigma^1 \subseteq \sigma^a$ we have:
    \begin{equation*}
        \pi^{\sigma^a}\Bigg(\Bigg[\bigwedge_{i \in A}f_i \wedge\bigwedge_{j\in B} e_i \Bigg].\Bigg[\bigwedge_{i=1}^{\mathrm{dim}F_{\sigma^a} -1}(v_i^{\sigma^1})^* \wedge\bigwedge_{i=\mathrm{rank}V+1}^{\mathrm{dim}F_{\sigma^a}}f_i^* \wedge\bigwedge_{j=1}^{n-\mathrm{dim}F_{\sigma^a}}e_j^* \Bigg] \Bigg) = 0.
    \end{equation*}

    Thus every coordinate of $\alpha$ contributes to zero to $\overline{\psi}(\alpha)$ except the ones for which $k=n-\mathrm{dim}F_{\sigma^a}$, which corresponds to $B=\{1,...,n-\mathrm{dim}F_{\sigma^a}\}$ Thus we obtain:
    \begin{equation*}
        \overline{\psi}(\alpha) = \overline{\psi} \Bigg( \sum_{A \in \mathcal{P}_{p-n+\mathrm{dim} F_{\sigma^a}}(\{1,...,\mathrm{dim}F_{\sigma^a}\})}\alpha_{A,\{1,...,n-\mathrm{dim} F_{\sigma^a}\} } \bigwedge_{i \in A}f_i \wedge\bigwedge_{j=1}^{n-\mathrm{dim} F_{\sigma^a}} e_i\Bigg).
    \end{equation*}

    Let us denote $\overline{\alpha}:=\sum_{A \in \mathcal{P}_{p-n+\mathrm{dim} F_{\sigma^a}}(\{1,...,\mathrm{rank}V\})}\alpha_{A,\{1,...,n-\mathrm{dim} F_{\sigma^a}\} } \bigwedge_{i \in A}f_i$. As we assume the triangulation to satisfy the hypothesis 1 for the integral domain $R$, we have for every $\sigma^1 \subseteq \sigma^a$ that $|\sigma^1|_r$ does not vanish in the integral domain $R$. Thus we have:
    \begin{align*}
        &\overline{\psi}(\alpha) = 0 \iff \forall \sigma^1 \subseteq \sigma^a, ~\overline{\alpha}.\Bigg[\bigwedge_{i=1}^{\mathrm{rank}V -1}(v_i^{\sigma^1})^*\wedge\bigwedge_{i=\mathrm{rank}V+1}^{\mathrm{dim}F_{\sigma^a}} f_i^*\Bigg] =0.
    \end{align*}

    Let us identify the vectors $(v_i^{\sigma^1})^*$ and the vector $u_{\sigma^1}^*$, for every $\sigma^1 \subseteq \sigma^a$, with their projection in $\frac{N_R}{V^{\perp}}=V^*$. Then for every $\sigma^1 \subseteq \sigma^a$ we have that $(u_{\sigma^1}^*,(v_1^{\sigma^1})^*,...,(v_{\mathrm{rank}V-1}^{\sigma^1})^*)$ forms a basis of $V^*$, dual to the basis $(u_{\sigma^1},v^{\sigma^1}_1,...,v^{\sigma^1}_{\mathrm{dim}F_{\sigma^a}})$ of $V$. Then the vector $\bigwedge_{i=1}^{\mathrm{rank}V -1}(v_i^{\sigma^1})^*$ is uniquely determined from the vector $u_{\sigma^1}$, up to an invertible factor of $R$. We have the following equality of rank:
    \begin{equation*}
        \mathrm{rank} ~ \sum_{\sigma^1 \subseteq \sigma^a} \bigwedge_{i=1}^{\mathrm{rank}V -1}(v_i^{\sigma^1})^* = \mathrm{rank}~ \sum_{\sigma^1 \subseteq \sigma^a} Ru_{\sigma^1} = \mathrm{rank}~ V = \mathrm{rank}~\bigwedge^{\mathrm{rank} V -1}V^*.
    \end{equation*}

    Thus we obtain that the family $( \bigwedge_{i=1}^{\mathrm{rank}V-1}(v_i^{\sigma^1})^*)_{\sigma^1 \subseteq \sigma^a}$ generates $\bigwedge^{\mathrm{rank} V -1}V^*$. Then we have:
    \begin{align*}
        &\overline{\psi}(\alpha) = 0 \\
        &\iff \forall x \in\bigwedge^{\mathrm{rank} V -1}\Bigg(\sum_{i=1}^{\mathrm{rank}V}Rf_i^*\Bigg), ~~\overline{\alpha}.\Bigg[x\wedge\bigwedge_{i=\mathrm{rank}V+1}^{\mathrm{dim}F_{\sigma^a}} f_i^*\Bigg] =0 \\
        & \iff \forall j \in \{1,...,\mathrm{rank}~V\},~~\overline{\alpha}.\Bigg[ \bigwedge_{i\in \{1,...,\mathrm{rank}V\}\backslash \{j\}} f_i^*\wedge\bigwedge_{i=\mathrm{rank}V+1}^{\mathrm{dim}F_{\sigma^a}} f_i^*\Bigg] =0 \\
        & \iff \forall A \in \mathcal{P}_{p-n+\mathrm{dim} F_{\sigma^a}}(\{1,...,\mathrm{dim}F_{\sigma^a}\}), ~~~~A \not\supseteq \{1,...,\mathrm{rank}V\} \implies\alpha_{A,\{1,...n-\mathrm{dim} F_{\sigma^a}\}}=0.
    \end{align*}

    Then we obtain:
    \begin{equation*}
        \mathrm{rank} ~\mathrm{ker} ~\overline{\psi}  = \mathrm{rank} ~ \bigwedge^p M_R - \Bigg[\binom{\mathrm{dim}F_{\sigma^a}}{p-n+\mathrm{dim}F_{\sigma^a}}-\binom{\mathrm{dim}F_{\sigma^a}-\mathrm{rank}V}{p-n+\mathrm{dim}F_{\sigma^a}-\mathrm{rank}V}\Bigg].
    \end{equation*}

    Finally we obtain:
    \begin{equation*}
        \mathrm{rank}~ \phi_{\sigma^{a},(p)} = \mathrm{rank} ~\overline{\psi} = \binom{\mathrm{dim}F_{\sigma^a}}{p-n+\mathrm{dim}F_{\sigma^a}}-\binom{\mathrm{dim}F_{\sigma^a}-\mathrm{rank}V}{p-n+\mathrm{dim}F_{\sigma^a}-\mathrm{rank}V}.
    \end{equation*}

    Now let us assume that the triangulation satisfies the hypothesis 2 for the integral domain $R$, and let us prove that the image of $\phi_{\sigma^{a},(p)}$ is saturated in $C_{a-1}(\mathcal{F}_{n-1-p}^R(*,\sigma^a))$. We just need to prove that the image of $\overline{\psi}$ is saturated. Let $y \in C_{a-1}(\mathcal{F}_{n-1-p}^R(*,\sigma^a))$ and $\mu \in R\backslash\{0\}$ such that $\overline{\psi}(\alpha) = \mu y$. For $\sigma^1 \subseteq \sigma^a$, using the properties of contraction we have:
    \begin{equation*}
        \mu y_{\sigma^1} =\pi^{\sigma^a} \Bigg( \alpha.(|\sigma^1|_rg(o(\sigma^1))\otimes1_R) \Bigg)= \pm\overline{\alpha}.\Bigg[\frac{|\sigma^1|_r}{\lambda_{\sigma^1}}\bigwedge_{i=1}^{\mathrm{rank}V -1}(v_i^{\sigma^1})^*\wedge\bigwedge_{i=\mathrm{rank}V+1}^{\mathrm{dim}F_{\sigma^a}} f_i^*\Bigg].
    \end{equation*}
    However we have already proven that the family  $( \bigwedge_{i=1}^{\mathrm{rank}V-1}(v_i^{\sigma^1})^*)_{\sigma^1 \subseteq \sigma^a}$ generates $\bigwedge^{n-1}\sum_{i=1}^{\mathrm{rank}V}Rf_i^*$. Then for any $1 \leq i\leq \mathrm{dim}V$ there exists a family of coefficients $(a^i_{\sigma^1})_{\sigma^1\subseteq \sigma^a}$ such that:
    \begin{equation*}
        \bigwedge_{j\in\{1,...,\mathrm{rank}V\}\backslash\{i\}}f_j^*=\sum_{\sigma^1\subseteq \sigma^a}a^i_{\sigma^1}\bigwedge_{i=1}^{\mathrm{rank}V-1}(v_i^{\sigma^1})^*.
    \end{equation*}
    Then for any $A \in \mathcal{P}_{p-n+\mathrm{dim} F_{\sigma^a}}(\{1,...,\mathrm{dim}F_{\sigma^a}\})$, if $A \not\supseteq\{1,...,\mathrm{rank}V\}$, there exists $i\in\{1,...,\mathrm{rank}V\}$ such that $i\not\in A$ and we have:
    \begin{equation*}
        \alpha_{A,\{1,...,n-\mathrm{dim}F_{\sigma^a}\}} \bigwedge_{j\in\{1,...\mathrm{rank}V\}\backslash(A\cup\{i\})}f_j^* = \sum_{\sigma^1\subseteq \sigma^a}\pm a^i_{\sigma^1}\frac{\lambda_{\sigma^1}}{|\sigma^1|_r} \mu y_{\sigma^1}.
    \end{equation*}
    We have used here the assumption that the triangulation satisfies the hypothesis $2$ for the integral domain $R$, which enable us to consider an inverse of $|\sigma^1|_r$ in $R$.
    
    From the last equation, $\mu$ divides $\alpha_{A,\{1,...,n-\mathrm{dim}F_{\sigma^a}\}}$ for any $A \in \mathcal{P}_{p-n+\mathrm{dim} F_{\sigma^a}}(\{1,...,\mathrm{dim}F_{\sigma^a}\})$ such that $A \not\supseteq\{1,...,\mathrm{rank}V\}$. However, we already know that if $A \supseteq\{1,...,\mathrm{rank}V\}$, then the term $\alpha_{A,\{1,...,n-\mathrm{dim}F_{\sigma^a}\}}$ contributes to zero in $\phi_{\sigma^{a},(p)}(\alpha)$. Then, for any $A \in \mathcal{P}_{p-n+\mathrm{dim} F_{\sigma^a}}(\{1,...,\mathrm{dim}F_{\sigma^a}\})$ such that $A \not\supseteq\{1,...,\mathrm{rank}V\}$ we can define $\beta_A$ the element of $R$ such that $\mu\beta_A = \alpha_{A,\{1,...,\mathrm{dim}F_{\sigma^a}\}}$ and we obtain:
    \begin{equation*}
        \phi_{\sigma^{a},(p)}\Bigg( \sum_{A \in \mathcal{P}_{p-n+\mathrm{dim} F_{\sigma^a}}(\{1,...,\mathrm{dim}F_{\sigma^a}\}), ~A\not\supseteq\{1,...\mathrm{rank}V\}}\beta_A\bigwedge_{i\in A}f_i^*\wedge \bigwedge_{j=1}^{n-\mathrm{dim}F_{\sigma^a}}e_j^* \Bigg) = y.
    \end{equation*}
    It proves that the image of $\phi_{\sigma^{a},(p)}$ is saturated in $C_{a-1}(\mathcal{F}_{n-1-p}^R(*,\sigma^a))$.
\end{proof}

\end{document}
